\documentclass[leqno, a4paper, 12pt]{jsarticle}
\usepackage{amssymb}
\usepackage{amsthm}
\usepackage{amsmath}
\usepackage{comment}
\usepackage{url}

\usepackage{enumitem}
%\usepackage{showkeys}


%定理環境 thmはあまり使わずpropをなるべく使う。
%主張は基本的に証明の中で使い、そのため番号を付けない。
%注意環境を付けてみた。
\newtheorem{thm}{定理}
\newtheorem{prop}[thm]{命題}
\newtheorem{cor}[thm]{系}
\newtheorem{lem}[thm]{補題}
\newtheorem{df}[thm]{定義}
\newtheorem{exam}[thm]{例}
\newtheorem{fact}[thm]{事実}
\newtheorem*{claim}{主張}
\newtheorem*{cau}{注意}
\newtheorem*{rmk}{注意}
\renewcommand{\proofname}{証明}
%矢印たち。
%\toは\rightarrowと同じ。\getsは\leftarrowと同じ。同値は\iff
\newcommand{\To}{\Longrightarrow}
\newcommand{\Gets}{\LongLeftarrow}
%黒板太字
\newcommand{\nn}{\mathbb{N}}
\newcommand{\zz}{\mathbb{Z}}
\newcommand{\qq}{\mathbb{Q}}
\newcommand{\rr}{\mathbb{R}}
\newcommand{\cc}{\mathbb{C}}
%無限大
\newcommand{\mgn}{\infty}
%差集合は\setminusで統一する。等号の否定は\neq
%真部分集合は\subsetneqq　偏微分は\partial
%ディスプレイスタイル。\dfracでディスプレイ分数
\newcommand{\dpst}{\displaystyle}

\title{論文メモ
\large{\\--- テイラー級数の収束と円周上の部分集合について ---}
}
\author{ナンブ　キトラ}
\date{\today}
\begin{document}
\maketitle

本棚を整理したいので，
溜まっている論文のメモを書いて未来への手紙とすることにする．


無限級数の収束に関する話である．

論文
\cite{MR0058734}でエルデシュは
間隙級数の冪級数で$|z|\le 1$では一様収束するが，
絶対収束はしないという級数を構成している．
エルデシュが初めてというわけではなく，
ハーディらのそういう研究を受けて発展として論文を書いたようだ．

論文\cite{MR0031049}と
\cite{MR0053210}
は連作である．
これらの論文では
$0$に収束する複素数列
$\{a_{n}\}$で
$\sum_{i=0}^{\infty}|a_{i}|=\infty$
となるものを係数とする複素関数
$\sum_{i=0}^{\infty}a_{i}z_{i}$
(テイラー級数)
を
考え，
この複素関数の単位円周
$C$
上での収束する点全体の集合を調べている．
このような級数は$C$上では絶対収束しないことに注意しよう．
テイラー級数の収束性とボレル階層が交わり面白い論文に見える．カントールの集合論がフーリエ級数の収束性に端を発していたことを思い返すと
この研究はある意味正統派の集合論とも言えるし先祖返りとも言えるかもしれない．

論文\cite{MR0031049}においては次の定理を証明している．
\begin{thm}
以下が成り立つ．
\begin{enumerate}[label=\textup{(\arabic*)}]
\item 
$M$を
$C$の
任意の
$F_{\sigma}$集合とすると
Taylor 級数
$\sum_{i=0}^{\infty}a_{i}z_{i}$が存在して
$M$上では収束し，
$C\setminus M$
上では
発散する．

\item 
$F$を
$C$の任意の閉集合とすると
テイラー級数が存在してそれは
$F$上では一様収束し，
$C\setminus F$
上では発散する．

\item 
$M$を
$C$の部分集合であって，
あるテイラー級数が$M$上一様収束し，
$C\setminus M$
が発散するようなものとする．
このとき
$M$は
$C$の閉集合である．

\item 

$M$
を
$C$の部分集合であって
あるテイラー級数が$M$上
収束し，
$C\setminus M$
では発散するものとする．
このとき
$M$は
$F_{\sigma\delta}$
集合である．

\end{enumerate}
\end{thm}
論文\cite{MR0053210}
は，正直何やってんのかよくわからないが
何やら単位開円盤上の複素関数とその境界上の収束に関する
微妙な話題の反例をたくさん与えているようである．

この論文には
schlicht function というのが出てくるがこれは単葉関数という意味である．
大まかにかいつまんで定理のいくつかを紹介する．
\begin{thm}
以下が成り立つ．
\begin{enumerate}
\item 
全てのテイラー級数が$F_{\sigma}$型とは限らない．
つまりテイラー級数であって，その円周上の収束する点全体が
$F_{\sigma}$に
は成らないものが存在する．
\item 
単位開円盤上で定義された有界単葉関数で，
$C$上では各点で収束するけれども，
$C$上のどの弧上でも一様収束には成らない
ようなものが存在する．

\item 
集合$|z|\le 1$上で連続な関数であって，
そのテイラー級数は$C$上で各点収束するけれども
どの弧の上でも一様収束しないようなものが存在する．

\item 
$C$の部分集合
$M$についてこれがあるテイラー級数の有界集合全体と一致するための必要十分条件は
$M$が閉集合で疎であることである．
\end{enumerate}
\end{thm}


論文
\cite{MR0043889}
では
$|z|\le 1$上で定義された
単葉関数
で円周上では
収束するけれども
絶対収束しない例を二つ構成している．
この論文は
パシフィックのVol.1なのでそれで珍しがって手元にあったのかもしれない．


論文
\cite{Mazurkiewicz1922}
はフランス語なのでよくわかっていないけど，
以下のようなことを示していると書いているような気がする．
\begin{thm}
円周$C$の閉部分集合$A$を任意に与える．
このとき以下が成り立つ．
\begin{enumerate}[label=\textup{(\arabic*)}]
\item 
テイラー級数が存在してそれは
$A$上の各点では収束するが，
$C\setminus A$上では
発散する．
\item 
テイラー級数が存在してそれは
$A$上では発散して
$C\setminus A$
上では収束する．

\end{enumerate}
\end{thm}
つまりは上の論文の
前身的結果ということである．


論文\cite{MR0050680}
はちょっとよくわかんなかったんだけど，
何やらスリットの入った領域に関してリーマンの写像定理の写像を具体的に書いているような気がする．


論文
\cite{MR0942366}と
\cite{MR0549529}
はKnopp's core theorem 
に関するものらしいが，
あんまりよくわかんない．
多分級数の総和法に関する話だと思うがわかんない．

今からは解析関数と空間充填曲線の話である．

論文
\cite{MR0015154}
では
$|z|\le 1$上で定義された連続関数で,
$|z|\le 1$では解析的で，
$|z|=1$の像は$\rr^{2}$の正方形を含むような
ものを初めて構成した．
論文
\cite{MR0062219}, 
\cite{MR0050016}
や
\cite{MR0133439}
もそういう感じ．
たぶん
数学書``Space-filling curves''の対応する章の論文を
印刷したんだと思う．

%READ!
%myplain is the same to amsplain style. 
\nocite{*}
\bibliographystyle{myplaindoidoi}
\bibliography{Taylar.bib}



\end{document}